\documentclass{article}
\usepackage{amsmath,hyperref}
\title{An effect-side authority cut for a persistent door command}
\begin{document}
\maketitle

\paragraph{Scope and claim.}
A session, an authorized incoming request, and an authorized recurring writer effect are different facts. This development check uses extracted, pinned Open-RMF door mode selection, session/phase test, target velocity and paired joint command writes. The request delay, gateway, per-gate generation, authoritative rotation event, hold, and velocity-following geometry belong to the fixture. In the pinned door plugin the request callback resolves a door by name and reduces the input to a mode; the update loop consumes that stored mode. No inspection of that path supplies a per-gate generation at the writer. This is not a ROS delivery, stock-gateway, hardware, independent transfer, or I1 acceptance claim.

\paragraph{Rule and exact minimum distinction.}
At integer writer update \(t\), select CLOSE for gate \(g\) only if a request is present at that gate, its authority generation equals the current one for \(g\), and \(t\ge h_g\), the first tick not protected by the hold. Use the same source writer to compute the velocity from current position and retained prior command. Rotation and incoming messages take effect before this test. For one named queued request, the Boolean of current gate-and-generation equality is necessary and sufficient for this single effect query given the time inputs. This Boolean is not a transferable certificate after a generation change; retain provenance or supply a new qualified equality. For fixed readiness \(r\) and hold expiry \(h\), a non-interrupted action of length \(d\) fits before \(D\) exactly when \(\max(r,h)+d\le D\), conditional on continuously valid authority and separate plant admissibility. The independent gate receives its usual writer updates, even during the affected hold. Lean verifies the pointwise laws and the tokenless obstruction. Neither read-only OPEN nor the chosen OPEN target promises instantaneous clearance.

\paragraph{Matched witnesses.}
A crossed generation at tick 2 while held until tick 3 gives two traces with the same delivered target, delivery tick 2, public observations through that tick, and physical command history up to expiry. At expiry the old authority requires suppressing CLOSE; the current authorized request requires applying it. A tokenless rule cannot implement both. In the modeled four-tick crossing starting at 3, the stale case dips only \(1/512\) from the initial open position; the fresh case dips \(11/512\) and violates the strict \(8/512\) tolerance. The fresh case is not offered the crossing at that cut. Clearing the queue before processing tick 2 does not eliminate an old delayed arrival at tick 2. An ingress-only check rejects that stale arrival but fails when a request accepted at tick 1 remains pending or latched past the rotation. In the five-prior-tick residual witness an OPEN target after rotation still writes \(-5/64\) at the rotation tick. Authorization suppresses a new CLOSE target, not the physical influence of its previous effects.

\paragraph{Comparison.}
There are 210 source-writer runs per arm, using one received CLOSE, one ordered generation rotation, two prior residuals, and the stated grid of hold and delivery times. An independently specified direct schedule is the reference for both gates' effective targets at every tick. Scoped effect-side checking matches 210/210. Hold alone matches 90, ingress check 180, a queue flush 120, and a global hold with otherwise correct effect-side authority 84. These are exact counts on this grid, not reliability estimates. The diagnostic failures are not the fair I1 competitor. Burrows's \href{https://static.usenix.org/events/osdi06/tech/full_papers/burrows/burrows_html/}{Chubby, section 2.4} already validates a resource-scoped generation at the receiving server. Effect-side application of that rule to each persistent write with the same hold and current authority ties by the pointwise predicate. No Chubby implementation ran here, and this establishes no advantage over published fencing. Both approaches still require the source-side connection that actually installs and checks the epoch.

\paragraph{Candidate interface and freeze gate.}
A supplier must expose the gate's allocation identity, effect-scoped grant, lease end and boundary ordering; a writer must trace the exact consumed request version, the applicable live grant, prior command state, actual write, and independent updates on one execution. Clear physical passage requires its own valid observation or bound and must not be synthesized from the fence acknowledgment. Compare the frozen joint rule to effect-side fencing plus an applicable native scheduler and to a same-information direct composition, with equal qualifying reads, update signals, producer authority and hold cost. For crash or restart, require durable/nonreused allocation identity, no alternate unguarded effect path, and evidence that queued/latched effects do not survive outside the validated span. These assumptions are not provided by the fixture. Do not choose a held-out case or call this the I1 freeze until the common mechanism, source supplier of each live premise, and matched positive envelope are pinned. The development comparison to replay next is read-only qualified evidence versus a real guard spanning reuse of a persistent or re-created writer, including a missed or reordered epoch update and a positive unrelated-service case.
\end{document}
